\documentclass[12pt]{article}
\usepackage[a4paper,margin=1in]{geometry}
\usepackage[utf8]{inputenc}
\usepackage[T1]{fontenc}
\usepackage{amsmath,amssymb}
\usepackage{textcomp}
\usepackage{parskip}
\begin{document}
\begin{center}
{\Large \textbf{2005 STEP III}}
\end{center}

\subsection*{Page 1}
\begin{flushleft}
STEP III, 2005 2\par
\bigskip\textbf{Section A: Pure Mathematics}\par
\[\pi\]
\bigskip\textbf{Question 1.} Show that sin A = cos B if and only if A = (4n + 1) B for some integer n.\par
2 ±\par
Show also that sin x cos x 6 \sqrt2 for all values of x and deduce that there are no solutions\par
±\par
\[to the equation \{sin (sin x) = c\{os (cos x).\]
\{ \{\par
Sketch, on the same axes, the graphs of y = sin (sin x) and y = cos (cos x). Sketch, not on the\par
\[previous axes, the graph of y = sin (2 sin x).\]
dy xy\par
\bigskip\textbf{Question 2.} Find the general solution of the differential equation = , where a = 0 , and show\par
\[dx - x^2 + a^2 6\]
2 2 2 2\par
that it can be written in the form y (x + a ) = c , where c is an arbitrary constant. Sketch\par
this curve.\par
d\par
2 2\par
Find an expression for (x + y ) and show that\par
dx\par
\[d^2 c^2 8c2x2\]
2 2\par
\[(x + y ) = 2 1 + .\]
\[dx2 \lambda - (x^2 + a^2)2 \mu (x^2 + a^2)3\]
\textbf{(i)} Show that, if 0 < c < a2, the points on the curve whose distance from the origin is least\par
are 0 , .\par
± a\par
\gamma \delta\par
\textbf{(ii)} If c > a2, determine the points on the curve whose distance from the origin is least.\par
\bigskip\textbf{Question 3.} Let f (x) = x2 + px + q and g (x) = x2 + rx + s . Find an expression for f (g (x)) and hence\par
find a necessary and sufficient condition on a, b and c for it to be possible to write the quartic\par
expression x4 + ax3 + bx2 + cx + d in the form f (g (x)), for some choice of values of p, q, r\par
and s.\par
Show further that this condition holds if and only if it is possible to write the quartic expression\par
x4 + ax3 + bx2 + cx + d in the form (x2 + vx + w)2 k, for some choice of values of v, w and k.\par
-\par
\[Find the roots of the quartic equation x^4 4x3 + 10x2 12x + 4 = 0 .\]
- -\par
\end{flushleft}
\subsection*{Page 2}
\begin{flushleft}
STEP III, 2005 3\par
\bigskip\textbf{Question 4.} The sequence un (n = 1, 2, . . .) satisfies the recurrence relation\par
un+1\par
\[un+2 = (kun un+1)\]
un -\par
where k is a constant.\par
\[If u 1 = a and u 2 = b , where a and b are non-zero and b = ka , prove by induction that\]
b\par
\[u 2n = u 2n-1\]
a\par
\gamma \delta\par
\[u 2n+1 = cu 2n\]
for n > 1, where c is a constant to be found in terms of k, a and b. Hence express u 2n and\par
u 2n-1 in terms of a, b, c and n.\par
Find conditions on a, b and k in the three cases:\par
\textbf{(i)} the sequence un is geometric;\par
\textbf{(ii)} the sequence un has period 2;\par
\textbf{(iii)} the sequence un has period 4.\par
\bigskip\textbf{Question 5.} Let P be the point on the curve y = ax2 + bx + c (where a is non-zero) at which the gradient\par
is m. Show that the equation of the tangent at P is\par
(m b)2\par
\[y mx = c - .\]
- - 4a\par
\[Show that the curves y = a x^2 + b x + c and y = a x^2 + b x + c (where a and a are non-\]
1 1 1 2 2 2 1 2\par
zero) have a common tangent with gradient m if and only if\par
\[(a a )m^2 + 2(a b a b )m + 4a a (c c ) + a b^2 a b^2 = 0 .\]
2 1 1 2 2 1 1 2 2 1 2 1 1 2\par
- - - -\par
\[Show that, in the case a = a , the two curves have exactly one common tangent if and only\]
1 2\par
if they touch each other. In the case a = a , find a necessary and sufficient condition for the\par
1 2\par
two curves to have exactly one common tangent.\par
\end{flushleft}
\subsection*{Page 3}
\begin{flushleft}
STEP III, 2005 4\par
\bigskip\textbf{Question 6.} In this question, you may use without proof the results\par
\[4 cosh 3 y 3 cosh y = cosh(3y) and arcosh y = ln y + y^2 1 .\]
- -\par
\gamma p \delta\par
[ Note: arcoshy is another notation for cosh -1 y ]\par
\[Show that the equation x^3 3a2x = 2a3 cosh T is satisfied by 2a cosh 1 T and hence that, if\]
- 3\par
b   1\par
c2 > b3 > 0, one of the roots of the equation x3 3bx = 2c is u+ , where u = c + \sqrtc2 b3 3 .\par
- u -\par
\gamma \delta\par
\[Show that the other two roots of the equation x^3 3bx = 2c are the roots of the quadratic\]
\[b b^2 -\]
2 2\par
\[equation x + u + x + u + b = 0 , and find these roots in terms of u, b and ω, where\]
\[u u^2 -\]
\[ω = 1 ( 1 + i\sqrt\{3\}).\]
2 -\par
\[Solve completely the equation x^3 6x = 6 .\]
-\par
1 m\par
\bigskip\textbf{Question 7.} Show that if du = F(u) + c , then dx = F(x m ) + c , where m = 0.\par
Z u f (u) Z x f (xm) 6\par
Find:\par
\textbf{(i)} dx ;\par
Z xn x\par
-\par
\textbf{(ii)} dx .\par
\[Z \sqrt\{xn\} + x^2\]
\bigskip\textbf{Question 8.} In this question, a and c are distinct non-zero complex numbers. The complex conjugate of\par
∗\par
any complex number z is denoted by z .\par
Show that\par
2 ∗ ∗ ∗ ∗\par
\[a c = aa + cc ac ca\]
| - | - -\par
and hence prove that the triangle OAC in the Argand diagram, whose vertices are represented\par
∗ ∗ ∗\par
\[by 0, a and c respectively, is right angled at A if and only if 2aa = ac + ca .\]
a\par
′\par
Points P and P in the Argand diagram are represented by the complex numbers ab and ,\par
b∗\par
where b is a non-zero complex number. A circle in the Argand diagram has centre C and\par
passes through the point A, and is such that OA is a tangent to the circle. Show that the\par
′\par
point P lies on the circle if and only if the point P lies on the circle.\par
a\par
Conversely, show that if the points represented by the complex numbers ab and , for some\par
b∗\par
∗\par
non-zero complex number b with bb = 1 , both lie on a circle centre C in the Argand diagram\par
which passes through A, then OA is a tangent to the circle.\par
\end{flushleft}
\subsection*{Page 4}
\begin{flushleft}
STEP III, 2005 5\par
\bigskip\textbf{Section B: Mechanics}\par
\bigskip\textbf{Question 9.} Two particles, A and B, move without friction along a horizontal line which is perpendicular to\par
a vertical wall. The coefficient of restitution between the two particles is e and the coefficient\par
\[of restitution between particle B and the wall is also e, where 0 < e < 1. The mass of\]
\[particle A is 4em (with m > 0), and the mass of particle B is (1 e)2m .\]
-\par
\[Initially, A is moving towards the wall with speed (1 e)v (where v > 0) and B is moving\]
-\par
away from the wall and towards A with speed 2ev. The two particles collide at a distance d\par
from the wall. Find the speeds of A and B after the collision.\par
When B strikes the wall, it rebounds along the same line. Show that a second collision will\par
take place, at a distance de from the wall.\par
Deduce that further collisions will take place. Find the distance from the wall at which the\par
nth collision takes place, and show that the times between successive collisions are equal.\par
\bigskip\textbf{Question 10.} Two thin discs, each of radius r and mass m, are held on a rough horizontal surface with their\par
\[\pimg\]
centres a distance 6r apart. A thin light elastic band, of natural length 2\pir and modulus ,\par
is wrapped once round the discs, its straight sections being parallel. The contact between the\par
elastic band and the discs is smooth. The coefficient of static friction between each disc and\par
the horizontal surface is µ, and each disc experiences a force due to friction equal to µmg\par
when it is sliding.\par
The discs are released simultaneously. If the discs collide, they rebound and a half of their\par
total kinetic energy is lost in the collision.\par
\textbf{(i)} Show that the discs start sliding, but come to rest before colliding, if and only if\par
\[2 < µ < 1.\]
\textbf{(ii)} Show that, if the discs collide at least once, their total kinetic energy just before the\par
first collision is 4 mgr(2 3µ).\par
3 -\par
\textbf{(iii)} Show that if 4 > µ2 > 5 the discs come to rest exactly once after the first collision.\par
9 27\par
\end{flushleft}
\subsection*{Page 5}
\begin{flushleft}
STEP III, 2005 6\par
\bigskip\textbf{Question 11.} A horizontal spindle rotates freely in a fixed bearing. Three light rods are each attached by\par
one end to the spindle so that they rotate in a vertical plane. A particle of mass m is fixed\par
to the other end of each of the three rods. The rods have lengths a, b and c, with a > b > c\par
and the angle between any pair of rods is 2 \pi. The angle between the rod of length a and the\par
vertical is θ, as shown in the diagram.\par
a\par
θ\par
\[2 \pi\]
c 3\par
\[2 \pi\]
3 b\par
Find an expression for the energy of the system and show that, if the system is in equilibrium,\par
then\par
\[(b c)\sqrt\{3\}\]
\[tan θ = - .\]
- 2a b c\par
- -\par
Deduce that there are exactly two equilibrium positions and determine which of the two\par
equilibrium positions is stable.\par
Show that, for the system to make complete revolutions, it must pass through its position of\par
stable equilibrium with an angular velocity of at least\par
4gR\par
,\par
\[r a^2 + b^2 + c^2\]
\[where 2R2 = (a b)2 + (b c)2 + (c a)2 .\]
- - -\par
\end{flushleft}
\subsection*{Page 6}
\begin{flushleft}
STEP III, 2005 7\par
\bigskip\textbf{Section C: Probability and Statistics}\par
12 Five independent timers time a runner as she runs four laps of a track. Four of the timers\par
measure the individual lap times, the results of the measurements being the random variables\par
T 1 to T 4, each of which has variance σ2 and expectation equal to the true time for the lap.\par
The fifth timer measures the total time for the race, the result of the measurement being the\par
random variable T which has variance σ2 and expectation equal to the true race time (which\par
is equal to the sum of the four true lap times).\par
Find a random variable X of the form aT + b(T 1 + T 2 + T 3 + T 4), where a and b are constants\par
independent of the true lap times, with the two properties:\par
(1) whatever the true lap times, the expectation of X is equal to the true race time;\par
(2) the variance of X is as small as possible.\par
Find also a random variable Y of the form cT + d(T 1 + T 2 + T 3 + T 4), where c and d are\par
constants independent of the true lap times, with the property that, whatever the true lap\par
times, the expectation of Y 2 is equal to σ2.\par
In one particular race, T takes the value 220 seconds and (T 1 + T 2 + T 3 + T 4) takes the value\par
220.5 seconds. Use the random variables X and Y to estimate an interval in which the true\par
race time lies.\par
\bigskip\textbf{Question 13.} A pack of cards consists of n + 1 cards, which are printed with the integers from 0 to n.\par
A game consists of drawing cards repeatedly at random from the pack until the card printed\par
with 0 is drawn, at which point the game ends. After each draw, the player receives \pounds 1 if the\par
card drawn shows any of the integers from 1 to w inclusive but receives nothing if the card\par
drawn shows any of the integers from w + 1 to n inclusive.\par
\textbf{(i)} In one version of the game, each card drawn is replaced immediately and randomly in\par
the pack. Explain clearly why the probability that the player wins a total of exactly \pounds 3\par
is equal to the probability of the following event occurring: out of the first four cards\par
drawn which show numbers in the range 0 to w, the numbers on the first three are\par
non-zero and the number on the fourth is zero. Hence show that the probability that\par
\[w^3\]
the player wins a total of exactly \pounds 3 is equal to .\par
(w + 1)4\par
Write down the probability that the player wins a total of exactly \pounds r and hence find\par
the expected total win.\par
\textbf{(ii)} In another version of the game, each card drawn is removed from the pack. Show that\par
the expected total win in this version is half of the expected total win in the other\par
version.\par
\end{flushleft}
\subsection*{Page 7}
\begin{flushleft}
STEP III, 2005 8\par
\bigskip\textbf{Question 14.} In this question, you may use the result\par
\infty m\par
t m! (n m)!\par
\[dt = - ,\]
\[Z^0 (t + k)n+2 (n + 1)! kn-m+1\]
\[where m and n are positive integers with n > m , and where k > 0 .\]
The random variable V has density function\par
C k\par
a+1\par
x\par
a\par
\[f (x) = (0 6 x < ) ,\]
(x + k)2a+2 \infty\par
(2a + 1)!\par
\[where a is a positive integer. Show that C = .\]
a! a!\par
Show, by means of a suitable substitution, that\par
v a \infty a\par
x u\par
\[dx = du\]
\[Z^0 (x + k)2a+2 Z k^2 (u + k)2a+2\]
v\par
and deduce that the median value of V is k. Find the expected value of V .\par
The random variable V represents the speed of a randomly chosen gas molecule. The time\par
s\par
taken for such a particle to travel a fixed distance s is given by the random variable T = .\par
V\par
Show that\par
\infty\par
C k\par
a+1\par
x\par
a\par
\[P(T < t) = dx ( )\]
Z s (x + k)2a+2 ∗\par
t\par
and hence find the density function of T . You may find it helpful to make the substitution\par
s\par
\[u = in the integral ( ).\]
x ∗\par
Hence show that the product of the median time and the median speed is equal to the dis-\par
tance s, but that the product of the expected time and the expected speed is greater than s.\par
\end{flushleft}
\end{document}
